\section{{\v C}ech cohomology and cohomology}
\label{section-cech-cohomology-cohomology}

\begin{lemma}
\label{lemma-injective-trivial-cech}
Let $X$ be a ringed space.
Let $\mathcal{U} : U = \bigcup_{i \in I} U_i$ be an open covering.
Let $\mathcal{I}$ be an injective $\mathcal{O}_X$-module.
Then
$$
\check{H}^p(\mathcal{U}, \mathcal{I}) =
\left\{
\begin{matrix}
\mathcal{I}(U) & \text{if} & p = 0 \\
0 & \text{if} & p > 0
\end{matrix}
\right.
$$
\end{lemma}

\begin{proof}
An injective $\mathcal{O}_X$-module is also injective as an object in
the category $\textit{PMod}(\mathcal{O}_X)$ (for example since
sheafification is an exact left adjoint to the inclusion functor,
using Homology, Lemma \ref{homology-lemma-adjoint-preserve-injectives}).
Hence we can apply Lemma \ref{lemma-cech-cohomology-derived-presheaves}
(or its proof) to see the result.
\end{proof}

\begin{lemma}
\label{lemma-cech-cohomology}
Let $X$ be a ringed space.
Let $\mathcal{U} : U = \bigcup_{i \in I} U_i$ be an open covering.
There is a transformation
$$
\check{\mathcal{C}}^\bullet(\mathcal{U}, -)
\longrightarrow
R\Gamma(U, -)
$$
of functors
$\textit{Mod}(\mathcal{O}_X) \to D^{+}(\mathcal{O}_X(U))$.
In particular this provides canonical maps
$\check{H}^p(\mathcal{U}, \mathcal{F}) \to H^p(U, \mathcal{F})$ for
$\mathcal{F}$ ranging over $\textit{Mod}(\mathcal{O}_X)$.
\end{lemma}

\begin{proof}
Let $\mathcal{F}$ be an $\mathcal{O}_X$-module. Choose an injective resolution
$\mathcal{F} \to \mathcal{I}^\bullet$. Consider the double complex
$\check{\mathcal{C}}^\bullet(\mathcal{U}, \mathcal{I}^\bullet)$ with terms
$\check{\mathcal{C}}^p(\mathcal{U}, \mathcal{I}^q)$.
There is a map of complexes
$$
\alpha :
\Gamma(U, \mathcal{I}^\bullet)
\longrightarrow
\text{Tot}(\check{\mathcal{C}}^\bullet(\mathcal{U}, \mathcal{I}^\bullet))
$$
coming from the maps
$\mathcal{I}^q(U) \to \check{H}^0(\mathcal{U}, \mathcal{I}^q)$
and a map of complexes
$$
\beta :
\check{\mathcal{C}}^\bullet(\mathcal{U}, \mathcal{F})
\longrightarrow
\text{Tot}(\check{\mathcal{C}}^\bullet(\mathcal{U}, \mathcal{I}^\bullet))
$$
coming from the map $\mathcal{F} \to \mathcal{I}^0$.
We can apply
Homology, Lemma \ref{homology-lemma-double-complex-gives-resolution}
to see that $\alpha$ is a quasi-isomorphism.
Namely, Lemma \ref{lemma-injective-trivial-cech} implies that
the $q$th row of the double complex
$\check{\mathcal{C}}^\bullet(\mathcal{U}, \mathcal{I}^\bullet)$ is a
resolution of $\Gamma(U, \mathcal{I}^q)$.
Hence $\alpha$ becomes invertible in $D^{+}(\mathcal{O}_X(U))$ and
the transformation of the lemma is the composition of $\beta$
followed by the inverse of $\alpha$. We omit the verification
that this is functorial.
\end{proof}

\begin{lemma}
\label{lemma-cech-h1}
Let $X$ be a topological space. Let $\mathcal{H}$ be an abelian sheaf
on $X$. Let $\mathcal{U} : X = \bigcup_{i \in I} U_i$ be an open covering.
The map
$$
\check{H}^1(\mathcal{U}, \mathcal{H}) \longrightarrow H^1(X, \mathcal{H})
$$
is injective and identifies $\check{H}^1(\mathcal{U}, \mathcal{H})$ via
the bijection of Lemma \ref{lemma-torsors-h1}
with the set of isomorphism classes of $\mathcal{H}$-torsors
which restrict to trivial torsors over each $U_i$.
\end{lemma}

\begin{proof}
To see this we construct an inverse map. Namely, let $\mathcal{F}$ be a
$\mathcal{H}$-torsor whose restriction to $U_i$ is trivial. By
Lemma \ref{lemma-trivial-torsor} this means there
exists a section $s_i \in \mathcal{F}(U_i)$. On $U_{i_0} \cap U_{i_1}$
there is a unique section $s_{i_0i_1}$ of $\mathcal{H}$ such that
$s_{i_0i_1} \cdot s_{i_0}|_{U_{i_0} \cap U_{i_1}} =
s_{i_1}|_{U_{i_0} \cap U_{i_1}}$. A computation shows
that $s_{i_0i_1}$ is a {\v C}ech cocycle and that its class is well
defined (i.e., does not depend on the choice of the sections $s_i$).
The inverse maps the isomorphism class of $\mathcal{F}$ to the cohomology
class of the cocycle $(s_{i_0i_1})$.
We omit the verification that this map is indeed an inverse.
\end{proof}

\begin{lemma}
\label{lemma-include}
Let $X$ be a ringed space.
Consider the functor
$i : \textit{Mod}(\mathcal{O}_X) \to \textit{PMod}(\mathcal{O}_X)$.
It is a left exact functor with right derived functors given by
$$
R^pi(\mathcal{F}) = \underline{H}^p(\mathcal{F}) :
U \longmapsto H^p(U, \mathcal{F})
$$
see discussion in Section \ref{section-locality}.
\end{lemma}

\begin{proof}
It is clear that $i$ is left exact.
Choose an injective resolution $\mathcal{F} \to \mathcal{I}^\bullet$.
By definition $R^pi$ is the $p$th cohomology {\it presheaf}
of the complex $\mathcal{I}^\bullet$. In other words, the
sections of $R^pi(\mathcal{F})$ over an open $U$ are given by
$$
\frac{\Ker(\mathcal{I}^p(U) \to \mathcal{I}^{p + 1}(U))}
{\Im(\mathcal{I}^{p - 1}(U) \to \mathcal{I}^p(U))}.
$$
which is the definition of $H^p(U, \mathcal{F})$.
\end{proof}

\begin{lemma}
\label{lemma-cech-spectral-sequence}
Let $X$ be a ringed space.
Let $\mathcal{U} : U = \bigcup_{i \in I} U_i$ be an open covering.
For any sheaf of $\mathcal{O}_X$-modules $\mathcal{F}$ there
is a spectral sequence $(E_r, d_r)_{r \geq 0}$ with
$$
E_2^{p, q} = \check{H}^p(\mathcal{U}, \underline{H}^q(\mathcal{F}))
$$
converging to $H^{p + q}(U, \mathcal{F})$.
This spectral sequence is functorial in $\mathcal{F}$.
\end{lemma}

\begin{proof}
This is a Grothendieck spectral sequence
(see
Derived Categories, Lemma \ref{derived-lemma-grothendieck-spectral-sequence})
for the functors
$$
i :  \textit{Mod}(\mathcal{O}_X) \to \textit{PMod}(\mathcal{O}_X)
\quad\text{and}\quad
\check{H}^0(\mathcal{U}, - ) : \textit{PMod}(\mathcal{O}_X)
\to \text{Mod}_{\mathcal{O}_X(U)}.
$$
Namely, we have $\check{H}^0(\mathcal{U}, i(\mathcal{F})) = \mathcal{F}(U)$
by Lemma \ref{lemma-cech-h0}. We have that $i(\mathcal{I})$ is
{\v C}ech acyclic by Lemma \ref{lemma-injective-trivial-cech}. And we
have that $\check{H}^p(\mathcal{U}, -) = R^p\check{H}^0(\mathcal{U}, -)$
as functors on $\textit{PMod}(\mathcal{O}_X)$
by Lemma \ref{lemma-cech-cohomology-derived-presheaves}.
Putting everything together gives the lemma.
\end{proof}

\begin{lemma}
\label{lemma-cech-spectral-sequence-application}
Let $X$ be a ringed space.
Let $\mathcal{U} : U = \bigcup_{i \in I} U_i$ be an open covering.
Let $\mathcal{F}$ be an $\mathcal{O}_X$-module.
Assume that $H^i(U_{i_0 \ldots i_p}, \mathcal{F}) = 0$
for all $i > 0$, all $p \geq 0$ and all $i_0, \ldots, i_p \in I$.
Then $\check{H}^p(\mathcal{U}, \mathcal{F}) = H^p(U, \mathcal{F})$
as $\mathcal{O}_X(U)$-modules.
\end{lemma}

\begin{proof}
We will use the spectral sequence of
Lemma \ref{lemma-cech-spectral-sequence}.
The assumptions mean that $E_2^{p, q} = 0$ for all $(p, q)$ with
$q \not = 0$. Hence the spectral sequence degenerates at $E_2$
and the result follows.
\end{proof}

\begin{lemma}
\label{lemma-ses-cech-h1}
Let $X$ be a ringed space.
Let
$$
0 \to \mathcal{F} \to \mathcal{G} \to \mathcal{H} \to 0
$$
be a short exact sequence of $\mathcal{O}_X$-modules.
Let $U \subset X$ be an open subset.
If there exists a cofinal system of open coverings $\mathcal{U}$
of $U$ such that $\check{H}^1(\mathcal{U}, \mathcal{F}) = 0$,
then the map $\mathcal{G}(U) \to \mathcal{H}(U)$ is
surjective.
\end{lemma}

\begin{proof}
Take an element $s \in \mathcal{H}(U)$. Choose an open covering
$\mathcal{U} : U = \bigcup_{i \in I} U_i$ such that
(a) $\check{H}^1(\mathcal{U}, \mathcal{F}) = 0$ and (b)
$s|_{U_i}$ is the image of a section $s_i \in \mathcal{G}(U_i)$.
Since we can certainly find $\mathcal{U}$ such that (b) holds
it follows from the assumptions of the lemma that we can find
$\mathcal{U}$ such that (a) and (b) both hold.
Consider the sections
$$
s_{i_0i_1} = s_{i_1}|_{U_{i_0i_1}} - s_{i_0}|_{U_{i_0i_1}}.
$$
Since $s_i$ lifts $s$ we see that $s_{i_0i_1} \in \mathcal{F}(U_{i_0i_1})$.
By the vanishing of $\check{H}^1(\mathcal{U}, \mathcal{F})$ we can
find sections $t_i \in \mathcal{F}(U_i)$ such that
$$
s_{i_0i_1} = t_{i_1}|_{U_{i_0i_1}} - t_{i_0}|_{U_{i_0i_1}}.
$$
Then clearly the sections $s_i - t_i$ satisfy the sheaf condition
and glue to a section of $\mathcal{G}$ over $U$ which maps to $s$.
Hence we win.
\end{proof}

\begin{lemma}
\label{lemma-cech-vanish}
\begin{slogan}
If higher {\v C}ech cohomology of an abelian sheaf vanishes for all open covers,
then higher cohomology vanishes.
\end{slogan}
Let $X$ be a ringed space.
Let $\mathcal{F}$ be an $\mathcal{O}_X$-module such that
$$
\check{H}^p(\mathcal{U}, \mathcal{F}) = 0
$$
for all $p > 0$ and any open covering $\mathcal{U} : U = \bigcup_{i \in I} U_i$
of an open of $X$. Then $H^p(U, \mathcal{F}) = 0$ for all $p > 0$
and any open $U \subset X$.
\end{lemma}

\begin{proof}
Let $\mathcal{F}$ be a sheaf satisfying the assumption of the lemma.
We will indicate this by saying ``$\mathcal{F}$ has vanishing higher
{\v C}ech cohomology for any open covering''.
Choose an embedding $\mathcal{F} \to \mathcal{I}$ into an
injective $\mathcal{O}_X$-module.
By Lemma \ref{lemma-injective-trivial-cech} $\mathcal{I}$ has vanishing higher
{\v C}ech cohomology for any open covering.
Let $\mathcal{Q} = \mathcal{I}/\mathcal{F}$
so that we have a short exact sequence
$$
0 \to \mathcal{F} \to \mathcal{I} \to \mathcal{Q} \to 0.
$$
By Lemma \ref{lemma-ses-cech-h1} and our assumptions
this sequence is actually exact as a sequence of presheaves!
In particular we have a long exact sequence of {\v C}ech cohomology
groups for any open covering $\mathcal{U}$, see
Lemma \ref{lemma-cech-cohomology-delta-functor-presheaves}
for example. This implies that $\mathcal{Q}$ is also an $\mathcal{O}_X$-module
with vanishing higher {\v C}ech cohomology for all open coverings.

\medskip\noindent
Next, we look at the long exact cohomology sequence
$$
\xymatrix{
0 \ar[r] &
H^0(U, \mathcal{F}) \ar[r] &
H^0(U, \mathcal{I}) \ar[r] &
H^0(U, \mathcal{Q}) \ar[lld] \\
&
H^1(U, \mathcal{F}) \ar[r] &
H^1(U, \mathcal{I}) \ar[r] &
H^1(U, \mathcal{Q}) \ar[lld] \\
&
\ldots & \ldots & \ldots \\
}
$$
for any open $U \subset X$. Since $\mathcal{I}$ is injective we
have $H^n(U, \mathcal{I}) = 0$ for $n > 0$ (see
Derived Categories, Lemma \ref{derived-lemma-higher-derived-functors}).
By the above we see that $H^0(U, \mathcal{I}) \to H^0(U, \mathcal{Q})$
is surjective and hence $H^1(U, \mathcal{F}) = 0$.
Since $\mathcal{F}$ was an arbitrary $\mathcal{O}_X$-module with
vanishing higher {\v C}ech cohomology we conclude that also
$H^1(U, \mathcal{Q}) = 0$ since $\mathcal{Q}$ is another of these
sheaves (see above). By the long exact sequence this in turn implies
that $H^2(U, \mathcal{F}) = 0$. And so on and so forth.
\end{proof}

\begin{lemma}
\label{lemma-cech-vanish-basis}
(Variant of Lemma \ref{lemma-cech-vanish}.)
Let $X$ be a ringed space.
Let $\mathcal{B}$ be a basis for the topology on $X$.
Let $\mathcal{F}$ be an $\mathcal{O}_X$-module.
Assume there exists a set of open coverings $\text{Cov}$
with the following properties:
\begin{enumerate}
\item For every $\mathcal{U} \in \text{Cov}$
with $\mathcal{U} : U = \bigcup_{i \in I} U_i$ we have
$U, U_i \in \mathcal{B}$ and every $U_{i_0 \ldots i_p} \in \mathcal{B}$.
\item For every $U \in \mathcal{B}$ the open coverings of $U$
occurring in $\text{Cov}$ is a cofinal system of open coverings
of $U$.
\item For every $\mathcal{U} \in \text{Cov}$ we have
$\check{H}^p(\mathcal{U}, \mathcal{F}) = 0$ for all $p > 0$.
\end{enumerate}
Then $H^p(U, \mathcal{F}) = 0$ for all $p > 0$ and any $U \in \mathcal{B}$.
\end{lemma}

\begin{proof}
Let $\mathcal{F}$ and $\text{Cov}$ be as in the lemma.
We will indicate this by saying ``$\mathcal{F}$ has vanishing higher
{\v C}ech cohomology for any $\mathcal{U} \in \text{Cov}$''.
Choose an embedding $\mathcal{F} \to \mathcal{I}$ into an
injective $\mathcal{O}_X$-module.
By Lemma \ref{lemma-injective-trivial-cech} $\mathcal{I}$
has vanishing higher {\v C}ech cohomology for any $\mathcal{U} \in \text{Cov}$.
Let $\mathcal{Q} = \mathcal{I}/\mathcal{F}$
so that we have a short exact sequence
$$
0 \to \mathcal{F} \to \mathcal{I} \to \mathcal{Q} \to 0.
$$
By Lemma \ref{lemma-ses-cech-h1} and our assumption (2)
this sequence gives rise to an exact sequence
$$
0 \to \mathcal{F}(U) \to \mathcal{I}(U) \to \mathcal{Q}(U) \to 0.
$$
for every $U \in \mathcal{B}$. Hence for any $\mathcal{U} \in \text{Cov}$
we get a short exact sequence of {\v C}ech complexes
$$
0 \to
\check{\mathcal{C}}^\bullet(\mathcal{U}, \mathcal{F}) \to
\check{\mathcal{C}}^\bullet(\mathcal{U}, \mathcal{I}) \to
\check{\mathcal{C}}^\bullet(\mathcal{U}, \mathcal{Q}) \to 0
$$
since each term in the {\v C}ech complex is made up out of a product of
values over elements of $\mathcal{B}$ by assumption (1).
In particular we have a long exact sequence of {\v C}ech cohomology
groups for any open covering $\mathcal{U} \in \text{Cov}$.
This implies that $\mathcal{Q}$ is also an $\mathcal{O}_X$-module
with vanishing higher {\v C}ech cohomology for all
$\mathcal{U} \in \text{Cov}$.

\medskip\noindent
Next, we look at the long exact cohomology sequence
$$
\xymatrix{
0 \ar[r] &
H^0(U, \mathcal{F}) \ar[r] &
H^0(U, \mathcal{I}) \ar[r] &
H^0(U, \mathcal{Q}) \ar[lld] \\
&
H^1(U, \mathcal{F}) \ar[r] &
H^1(U, \mathcal{I}) \ar[r] &
H^1(U, \mathcal{Q}) \ar[lld] \\
&
\ldots & \ldots & \ldots \\
}
$$
for any $U \in \mathcal{B}$. Since $\mathcal{I}$ is injective we
have $H^n(U, \mathcal{I}) = 0$ for $n > 0$ (see
Derived Categories, Lemma \ref{derived-lemma-higher-derived-functors}).
By the above we see that $H^0(U, \mathcal{I}) \to H^0(U, \mathcal{Q})$
is surjective and hence $H^1(U, \mathcal{F}) = 0$.
Since $\mathcal{F}$ was an arbitrary $\mathcal{O}_X$-module with
vanishing higher {\v C}ech cohomology for all $\mathcal{U} \in \text{Cov}$
we conclude that also $H^1(U, \mathcal{Q}) = 0$ since $\mathcal{Q}$ is
another of these sheaves (see above). By the long exact sequence this in
turn implies that $H^2(U, \mathcal{F}) = 0$. And so on and so forth.
\end{proof}

\begin{lemma}
\label{lemma-pushforward-injective}
Let $f : X \to Y$ be a morphism of ringed spaces.
Let $\mathcal{I}$ be an injective $\mathcal{O}_X$-module.
Then
\begin{enumerate}
\item $\check{H}^p(\mathcal{V}, f_*\mathcal{I}) = 0$
for all $p > 0$ and any open covering
$\mathcal{V} : V = \bigcup_{j \in J} V_j$ of $Y$.
\item $H^p(V, f_*\mathcal{I}) = 0$ for all $p > 0$ and
every open $V \subset Y$.
\end{enumerate}
In other words, $f_*\mathcal{I}$ is right acyclic for $\Gamma(V, -)$
(see
Derived Categories, Definition \ref{derived-definition-derived-functor})
for any $V \subset Y$ open.
\end{lemma}

\begin{proof}
Set $\mathcal{U} : f^{-1}(V) = \bigcup_{j \in J} f^{-1}(V_j)$.
It is an open covering of $X$ and
$$
\check{\mathcal{C}}^\bullet(\mathcal{V}, f_*\mathcal{I}) =
\check{\mathcal{C}}^\bullet(\mathcal{U}, \mathcal{I}).
$$
This is true because
$$
f_*\mathcal{I}(V_{j_0 \ldots j_p})
= \mathcal{I}(f^{-1}(V_{j_0 \ldots j_p})) =
\mathcal{I}(f^{-1}(V_{j_0}) \cap \ldots \cap f^{-1}(V_{j_p}))
= \mathcal{I}(U_{j_0 \ldots j_p}).
$$
Thus the first statement of the lemma follows from
Lemma \ref{lemma-injective-trivial-cech}. The second statement
follows from the first and Lemma \ref{lemma-cech-vanish}.
\end{proof}

\noindent
The following lemma implies in particular that
$f_* : \textit{Ab}(X) \to \textit{Ab}(Y)$ transforms injective
abelian sheaves into injective abelian sheaves.

\begin{lemma}
\label{lemma-pushforward-injective-flat}
Let $f : X \to Y$ be a morphism of ringed spaces.
Assume $f$ is flat.
Then $f_*\mathcal{I}$ is an injective $\mathcal{O}_Y$-module
for any injective $\mathcal{O}_X$-module $\mathcal{I}$.
\end{lemma}

\begin{proof}
In this case the functor $f^*$ transforms injections into injections
(Modules, Lemma \ref{modules-lemma-pullback-flat}).
Hence the result follows from
Homology, Lemma \ref{homology-lemma-adjoint-preserve-injectives}.
\end{proof}

\begin{lemma}
\label{lemma-cohomology-products}
Let $(X, \mathcal{O}_X)$ be a ringed space. Let $I$ be a set.
For $i \in I$ let  $\mathcal{F}_i$ be an $\mathcal{O}_X$-module.
Let $U \subset X$ be open. The canonical map
$$
H^p(U, \prod\nolimits_{i \in I} \mathcal{F}_i)
\longrightarrow
\prod\nolimits_{i \in I} H^p(U, \mathcal{F}_i)
$$
is an isomorphism for $p = 0$ and injective for $p = 1$.
\end{lemma}

\begin{proof}
The statement for $p = 0$ is true because the product of sheaves
is equal to the product of the underlying presheaves, see
Sheaves, Section \ref{sheaves-section-limits-sheaves}.
Proof for $p = 1$. Set $\mathcal{F} = \prod \mathcal{F}_i$.
Let $\xi \in H^1(U, \mathcal{F})$ map to zero in
$\prod H^1(U, \mathcal{F}_i)$. By locality of cohomology, see
Lemma \ref{lemma-kill-cohomology-class-on-covering},
there exists an open covering $\mathcal{U} : U = \bigcup U_j$ such that
$\xi|_{U_j} = 0$ for all $j$. By Lemma \ref{lemma-cech-h1} this means
$\xi$ comes from an element
$\check \xi \in \check H^1(\mathcal{U}, \mathcal{F})$.
Since the maps
$\check H^1(\mathcal{U}, \mathcal{F}_i) \to H^1(U, \mathcal{F}_i)$
are injective for all $i$ (by Lemma \ref{lemma-cech-h1}), and since
the image of $\xi$ is zero in $\prod H^1(U, \mathcal{F}_i)$ we see
that the image
$\check \xi_i = 0$ in $\check H^1(\mathcal{U}, \mathcal{F}_i)$.
However, since $\mathcal{F} = \prod \mathcal{F}_i$ we see that
$\check{\mathcal{C}}^\bullet(\mathcal{U}, \mathcal{F})$ is the
product of the complexes
$\check{\mathcal{C}}^\bullet(\mathcal{U}, \mathcal{F}_i)$,
hence by
Homology, Lemma \ref{homology-lemma-product-abelian-groups-exact}
we conclude that $\check \xi = 0$ as desired.
\end{proof}







